\section{Ex1: Phân lớp điểm dữ liệu bằng thuật toán KNN}

\textbf{Đề bài:} Given a collection dataset as following: Apply KNN to classify the data point $z = 5.0$ according to its $1$-, $3$-, $5$-, and $9$-nearest neighbors (using majority vote).

\textbf{Dữ liệu:} Tập hợp dữ liệu có các cặp giá trị:
\begin{itemize}
    \item $x$: $0.5, 3.0, 4.5, 4.6, 4.9, 5.2, 5.3, 5.5, 7.0, 9.5$
    \item $y$: $-, -, +, +, +, -, -, -, -, -$
\end{itemize}
Điểm cần phân loại: $z = 5.0$.

\subsection*{Bước 1: Tính khoảng cách Euclidean}
Khoảng cách được tính theo công thức:
$$d(x, z) = \sqrt{(x - z)^2}$$
Thay $z = 5.0$ vào, ta có bảng khoảng cách giữa $z$ và từng điểm $x_i$:

\begin{table}[H]
    \centering
    \begin{tabular}{|c|c|c|}
        \hline
        \textbf{$x$} & \textbf{$y$} & \textbf{Khoảng cách} \\
        \hline
        0.5 & $-$ & 4.5 \\ \hline
        3.0 & $-$ & 2.0 \\ \hline
        4.5 & $+$ & 0.5 \\ \hline
        4.6 & $+$ & 0.4 \\ \hline
        4.9 & $+$ & 0.1 \\ \hline
        5.2 & $-$ & 0.2 \\ \hline
        5.3 & $-$ & 0.3 \\ \hline
        5.5 & $-$ & 0.5 \\ \hline
        7.0 & $-$ & 2.0 \\ \hline
        9.5 & $-$ & 4.5 \\ \hline
    \end{tabular}
\end{table}

\subsection*{Bước 2: Sắp xếp khoảng cách}
Sắp xếp các điểm theo thứ tự tăng dần của khoảng cách:

\begin{table}[H]
    \centering
    \begin{tabular}{|c|c|c|}
        \hline
        \textbf{$x$} & \textbf{$y$} & \textbf{Khoảng cách} \\
        \hline
        4.9 & $+$ & 0.1 \\ \hline
        5.2 & $-$ & 0.2 \\ \hline
        5.3 & $-$ & 0.3 \\ \hline
        4.6 & $+$ & 0.4 \\ \hline
        4.5 & $+$ & 0.5 \\ \hline
        5.5 & $-$ & 0.5 \\ \hline
        3.0 & $-$ & 2.0 \\ \hline
        7.0 & $-$ & 2.0 \\ \hline
        0.5 & $-$ & 4.5 \\ \hline
        9.5 & $-$ & 4.5 \\ \hline
    \end{tabular}
\end{table}

\subsection*{Bước 3: Xác định nhãn theo $k$-láng giềng}
\begin{itemize}
    \item \textbf{Với $k = 1$:} Láng giềng gần nhất là $4.9$, nhãn $(+)$. Do đó, nhãn của $z$ là \textbf{$+$}.
    \item \textbf{Với $k = 3$:} 3 láng giềng gần nhất là $\{4.9 (+), 5.2 (-), 5.3 (-)\}$. Bầu phiếu đa số: $(+)$ có 1 phiếu, $(-)$ có 2 phiếu. Do đó, nhãn của $z$ là \textbf{$-$}.
    \item \textbf{Với $k = 5$:} 5 láng giềng gần nhất là $\{4.9 (+), 5.2 (-), 5.3 (-), 4.6 (+), 4.5 (+)\}$. Bầu phiếu đa số: $(+)$ có 3 phiếu, $(-)$ có 2 phiếu. Do đó, nhãn của $z$ là \textbf{$+$}.
    \item \textbf{Với $k = 9$:} 9 láng giềng gần nhất có 3 phiếu $(+)$ và 6 phiếu $(-)$. Do đó, nhãn của $z$ là \textbf{$-$}.
\end{itemize}